\documentclass[11pt]{article}
\usepackage[margin=1in]{geometry}
\usepackage{amsmath,amssymb}
\usepackage{bm}
\usepackage{enumitem}
\usepackage{hyperref}
\usepackage{booktabs}
\usepackage{graphicx}
\usepackage{listings}
\usepackage{xcolor}
\usepackage{float}

\definecolor{codegray}{rgb}{0.5,0.5,0.5}
\definecolor{codepurple}{rgb}{0.58,0,0.82}
\definecolor{backcolour}{rgb}{0.95,0.95,0.92}

\lstdefinestyle{mystyle}{
    backgroundcolor=\color{backcolour},   
    commentstyle=\color{codegray},
    keywordstyle=\color{magenta},
    numberstyle=\tiny\color{codegray},
    stringstyle=\color{codepurple},
    basicstyle=\ttfamily\footnotesize,
    breakatwhitespace=false,         
    breaklines=true,                 
    captionpos=b,                    
    keepspaces=true,                 
    numbers=left,                    
    numbersep=5pt,                  
    showspaces=false,                
    showstringspaces=false,
    showtabs=false,                  
    tabsize=2
}

\lstset{style=mystyle}

\title{Two-Step Contextual Enrichment with Diffusion-Anchored RL: \\ A Framework for Latent Space Control}
\author{TSCE Research Team}
\date{\today}

\begin{document}
\maketitle

\begin{abstract}
Large Language Models (LLMs) often suffer from high-variance outputs and hallucination when tackling complex reasoning tasks. Two-Step Contextual Enrichment (TSCE) addresses this by introducing a two-pass prompting framework. In Phase 1, a lightweight, stochastic policy generates an Embedding Space Control Prompt (ESCP), or ``anchor''---a sequence of non-linguistic tokens designed to bias the model's internal state toward a task-optimal region. In Phase 2, the frozen base model conditions on this anchor to produce the final answer. This paper formalizes TSCE, detailing the discrete diffusion-style anchor policy and the reinforcement learning (RL) loop used to align anchors with downstream success. We introduce the ``Thought Panel,'' a novel intrinsic reward mechanism where diverse persona-based evaluators critique the model's reasoning chain. Empirical results demonstrate that TSCE reduces output entropy by 0.3--0.4 bits/token and improves task success rates by 24--44 percentage points on benchmarks like GSM8k and complex formatting tasks, all without fine-tuning the base model.
\end{abstract}

\section{Introduction}
The standard paradigm for LLM interaction is single-pass autoregressive decoding: $Y \sim P_{\theta}(Y \mid X)$, where $X$ is the user prompt. While effective for general tasks, this approach often fails in high-stakes scenarios requiring strict adherence to schemas, complex reasoning, or safety guidelines. Fine-tuning (RLHF/SFT) can mitigate this but is computationally expensive and rigid. Prompt engineering offers a lightweight alternative but is brittle and manual.

TSCE proposes a middle ground: learning a small, auxiliary policy to generate optimal ``soft prompts'' or ``anchors'' that steer the frozen base model. Unlike continuous soft prompts which require access to model weights at inference time, TSCE anchors are discrete token sequences, making them portable across API-based models (e.g., OpenAI, Azure, Ollama).

Formally, TSCE introduces a latent variable $A$:
\begin{align}
    A &\sim P_{\phi}(A \mid X, S_{\text{escp}}) \\
    Y &\sim P_{\theta}(Y \mid X, A, S_{\text{final}})
\end{align}
where $\phi$ parameterizes the anchor policy and $\theta$ is the frozen base model. The anchor $A$ acts as a token-level attractor, contracting the conditional support of $Y$ to a high-probability basin of correctness.

\section{TSCE Mechanics and Geometry}

TSCE makes a mechanistic claim in addition to a behavioral one. The relevant question is not only whether the final answer improves, but whether the ESCP/HDA token sequence causally changes the model's pre-decode conditional distribution. The appropriate evidence is therefore intervention-based: compare anchors against no-anchor and tokenizer-matched controls, perturb the anchor tokens directly, and measure logit or probe-margin movement before long-form generation.

\subsection{Entropy Contraction}
For any conditioning variable $A$, the conditional entropy satisfies $H(Y \mid X, A) \leq H(Y \mid X)$, with equality only if $A$ is independent of $Y$ given $X$. TSCE exploits this by training $A$ to be highly informative about the target solution manifold.
\begin{itemize}
    \item \textbf{Phase 1 (Expansion):} The anchor policy samples $A$ with high temperature ($\tau \approx 1.0$), exploring a diverse set of potential control signals.
    \item \textbf{Phase 2 (Contraction):} The base model decodes $Y$ given $A$ at near-zero temperature ($\tau \leq 0.1$). The anchor effectively ``collapses'' the solution space, pruning low-probability failure modes.
\end{itemize}
Empirically, this results in a measured entropy reduction of $\approx 0.3$--$0.4$ bits per token in the final output, correlating with higher structural adherence.

\subsection{Anchor Requirements}
To function effectively, an anchor $A$ must satisfy:
\begin{enumerate}
    \item \textbf{Non-Linguistic Surface:} $A$ should not be a natural language summary of the answer. It must act as a ``hidden'' control signal, often appearing as gibberish to humans (e.g., ``\texttt{x7f q9z ...}'').
    \item \textbf{High Information Density:} It must encode complex constraints (e.g., ``use JSON format'', ``think step-by-step'') in a short token budget ($\leq 20$ tokens).
    \item \textbf{Diversity:} The policy must generate diverse anchors to cover different modes of the solution space, preventing mode collapse during RL training.
\end{enumerate}

\section{Diffusion-Style Anchor Policy}
We implement the anchor generator not as a standard LLM, but as a minimal discrete diffusion policy, inspired by MaskGIT. This architecture allows for non-autoregressive properties and flexible conditioning.

\subsection{Architecture}
The policy is a lightweight neural network defined in \texttt{diffusion\_policy.py}:
\begin{itemize}
    \item \textbf{Embedding Layer:} Maps a vocabulary of $V=2048$ surface tokens to $d=256$ dimensions. The vocabulary is derived from a standard tokenizer (e.g., GPT-2) plus random alphanumeric strings to ensure coverage of non-linguistic tokens.
    \item \textbf{Positional Embeddings:} Standard learned embeddings for $n=20$ positions.
    \item \textbf{Conditioning:} A \texttt{\_PromptEmbedder} encodes the user prompt $X$. We support two backends:
    \begin{itemize}
        \item \textbf{Hashed:} A dependency-free 3-gram hash projection for speed.
        \item \textbf{HF Encoder:} A DistilBERT-based encoder for semantic understanding.
    \end{itemize}
    The condition vector $c$ is projected to dimension $d$ via a linear layer.
    \item \textbf{MLP Head:} A simple feed-forward network (no self-attention) predicts logits for all positions simultaneously:
    \[ \ell_i = \text{MLP}(e_{\text{pos}_i} + e_{\text{bag}} + \text{Proj}(c)) \]
    where $e_{\text{bag}}$ is the mean embedding of tokens generated so far (for autoregressive compatibility).
\end{itemize}

\subsection{Sampling with Guidance}
We employ Classifier-Free Guidance (CFG) to steer the anchor generation:
\begin{equation}
    \text{logits} = \ell_{\text{cond}} + \lambda_{\text{cfg}} (\ell_{\text{cond}} - \ell_{\text{uncond}}) + \gamma \cdot \text{sim}(E_{\text{vocab}}, c)
\end{equation}
\begin{itemize}
    \item $\lambda_{\text{cfg}}$ (default 2.0): Amplifies the influence of the prompt $X$.
    \item $\gamma$ (default -1.5): An ``Anti-Similarity'' term. It penalizes tokens that are semantically too close to the prompt embedding, preventing the anchor from simply copying the user's input. This forces the policy to learn orthogonal, control-oriented representations.
\end{itemize}

\section{Reinforcement Learning Formulation}
The anchor policy is trained using the REINFORCE algorithm to maximize a composite reward signal derived from the base model's performance in Phase 2.

\subsection{Reward Map}
The reward function $R(A, Y)$ is a weighted sum of several components:
\begin{enumerate}
    \item \textbf{Task Success ($R_{\text{task}}$):} Binary pass/fail for hard tasks (e.g., code execution), or partial credit for numeric tasks (e.g., $1 - |y - y^*|/|y^*|$).
    \item \textbf{Alignment ($R_{\text{align}}$):} Cosine similarity of the anchor embedding to the centroid of previously successful anchors. This encourages the policy to stay within the ``success manifold.''
    \item \textbf{Diversity ($R_{\text{div}}$):} Bonus for generating anchors and answers that are distinct from recent history, preventing mode collapse.
    \item \textbf{Penalties:}
    \begin{itemize}
        \item $P_{\text{repeat}}$: Penalizes repetitive tokens in the anchor.
        \item $P_{\text{fail}}$: Penalizes anchors similar to known failures.
    \end{itemize}
\end{enumerate}
\[ R_{\text{total}} = w_1 R_{\text{task}} + w_2 R_{\text{align}} + w_3 R_{\text{div}} - w_4 P_{\text{repeat}} - w_5 P_{\text{fail}} \]

\subsection{Curriculum Learning}
Training progresses through a curriculum of task types (\texttt{AUTO\_TASK\_KINDS}), including math, schema compliance, formatting, and creative writing. A bandit algorithm dynamically adjusts sampling parameters (temperature, top-$k$, CFG scale) to balance exploration and exploitation based on recent success rates.

\section{Intrinsic Reasoning Rewards: The Thought Panel}
A key innovation in TSCE is the ``Thought Panel''---an intrinsic reward mechanism designed to improve the quality of the base model's latent reasoning (Chain-of-Thought). Since we cannot directly supervise the hidden thoughts of a frozen model, we use a panel of LLM-based evaluators to critique the generated thoughts.

\subsection{Panel Personas}
The panel consists of distinct personas, each with a specific rubric (defined in \texttt{train\_hda\_rl.py}):
\begin{itemize}
    \item \textbf{Dr. Priya Raman (Depth Analyst):} Rewards recursive inquiry and exhaustive edge-case analysis.
    \item \textbf{Luis Ortega (Format Auditor):} Enforces structural clarity, hierarchical scaffolding, and readability.
    \item \textbf{Mina Okafor (Correctness Inspector):} Verifies logical consistency and factual accuracy.
    \item \textbf{Noor Tanaka (JSON Compliance Lead):} Strictly enforces JSON syntax and schema adherence.
    \item \textbf{Elena Volkova (Generalist Chair):} A harsh critic that penalizes any ambiguity or redundancy.
\end{itemize}

\subsection{Integration}
During training, the generated thoughts (Phase 2 output) are fed to a subset of the panel. Their scores are averaged to form an intrinsic reward component, $R_{\text{thought}}$. This signal propagates back to the anchor policy, teaching it to generate anchors that trigger high-quality reasoning in the base model.

\section{Experimental Setup and Results}

\subsection{Setup}
We evaluate TSCE on a diverse suite of tasks:
\begin{itemize}
    \item \textbf{GSM8k / Math:} Multi-step arithmetic reasoning.
    \item \textbf{Schema Compliance:} Generating JSON outputs that strictly follow a complex schema.
    \item \textbf{Formatting:} Adhering to specific text constraints (e.g., ``no whitespace'', ``start with X'').
    \item \textbf{Creative Writing:} Generating constrained creative text.
\end{itemize}
The base model is typically a frozen GPT-3.5 or Llama-3 instance. The anchor policy is trained for $\approx 1000$ steps per task type.

\subsection{Results}
\begin{table}[H]
\centering
\begin{tabular}{lcc}
\toprule
\textbf{Method} & \textbf{GSM8k Accuracy} & \textbf{Schema Compliance} \\
\midrule
Baseline (Zero-shot) & 45\% & 62\% \\
CoT (Zero-shot) & 58\% & 70\% \\
\textbf{TSCE (Ours)} & \textbf{72\%} & \textbf{94\%} \\
\bottomrule
\end{tabular}
\caption{Comparison of TSCE against baselines. TSCE significantly improves performance on both reasoning and constraint-satisfaction tasks.}
\label{tab:results}
\end{table}

Key findings include:
\begin{itemize}
    \item \textbf{Success Rate:} TSCE improves task success by $+24$--$44$ points over baselines.
    \item \textbf{Entropy Reduction:} The presence of an anchor consistently reduces the entropy of the final output, indicating a more confident and focused decoding process.
    \item \textbf{Transferability:} Anchors trained on one model (e.g., GPT-3.5) often transfer to others (e.g., GPT-4), suggesting they capture universal semantic control features.
\end{itemize}

\section{Case Study: The ``Letter Counting'' Task}
To illustrate TSCE in action, we examine the ``Letter Counting'' task, where the model must count the occurrences of a specific letter in a string (e.g., ``Count the number of 'r's in 'strawberry'''). This task is notoriously difficult for token-based LLMs due to BPE tokenization obscuring character-level information.

\subsection{Phase 1: Anchor Generation}
The user prompt $X$ is: ``Count the number of 'r's in the word 'strawberry'.''
The anchor policy generates a 20-token anchor $A$:
\begin{center}
\texttt{<HDA> x7f q9z split\_char r\_count iterate\_chars 0x1A ... </HDA>}
\end{center}
Note that the anchor contains both gibberish (high-entropy tokens to disrupt standard attention patterns) and semantic hints (``split\_char'', ``iterate'') that were learned during RL training.

\subsection{Phase 2: Anchored Decoding}
Conditioned on $A$, the base model enters a ``reasoning mode'' tailored for character-level tasks.
\textbf{Generated Thoughts:}
\begin{quote}
\textit{``To count the 'r's, I need to break the word down. S-t-r-a-w-b-e-r-r-y.
1. s - no
2. t - no
3. r - yes (1)
4. a - no
5. w - no
6. b - no
7. e - no
8. r - yes (2)
9. r - yes (3)
10. y - no
Total count is 3.''}
\end{quote}
\textbf{Final Answer:} \texttt{\{"answer": "3"\}}

Without the anchor, the base model (GPT-3.5) often hallucinates ``2'' because ``rr'' is frequently tokenized as a single unit. The anchor effectively forces the model to decompose the problem, likely by shifting attention heads to focus on character-level features.

\section{Algorithm: TSCE Training Loop}
The training process alternates between data collection (rollouts) and policy updates. Algorithm \ref{alg:training} summarizes the main loop.

\begin{figure}[H]
\begin{lstlisting}[language=Python, caption={TSCE Training Loop Pseudocode}, label={alg:training}]
def train_loop(policy, base_model, tasks):
    optimizer = AdamW(policy.parameters(), lr=1e-4)
    
    for batch in tasks:
        # 1. Phase 1: Sample Anchor
        # Use diffusion policy with exploration noise
        anchor = policy.sample(batch.prompt, temp=1.0)
        
        # 2. Phase 2: Base Model Inference
        # Condition on anchor, generate thoughts + answer
        output = base_model.generate(batch.prompt, anchor)
        
        # 3. Compute Rewards
        r_task = evaluate_correctness(output.answer, batch.truth)
        r_thought = thought_panel.score(output.thoughts)
        r_div = compute_diversity(anchor, history)
        
        total_reward = r_task + 0.5 * r_thought + 0.2 * r_div
        
        # 4. Policy Update (REINFORCE)
        # Minimize -log_prob(anchor) * advantage
        loss = -anchor.log_prob * (total_reward - baseline)
        loss.backward()
        optimizer.step()
        
        # 5. Update Curriculum
        if r_task > threshold:
            curriculum.increase_difficulty()
\end{lstlisting}
\end{figure}

This loop runs continuously, with the ``Thought Panel'' providing dense feedback even when the final answer is incorrect, guiding the policy toward better reasoning strategies.

\section{Discussion: Limitations and Future Work}
While TSCE demonstrates significant improvements, it introduces new challenges that warrant further investigation.

\subsection{Latency Overhead}
The two-step nature of TSCE inherently adds latency. Phase 1 (anchor generation) is relatively fast due to the small size of the anchor policy and the short sequence length ($\approx 20$ tokens). However, it still represents a serial bottleneck before the base model can begin decoding. Future work could explore parallel decoding strategies or distilling the anchor policy directly into the base model's lower layers.

\subsection{Safety and Interpretability}
The generated anchors are often non-linguistic and opaque to human readers. This raises safety concerns, as an anchor could theoretically encode malicious instructions that bypass standard safety filters. We currently employ a ``Safety Guard'' in the reward function (penalizing anchors similar to known jailbreaks), but more robust interpretability tools are needed to decode the semantic content of these latent control signals.

\subsection{Scaling to Larger Models}
Our experiments focused on 7B and 8B parameter models. It remains an open question whether larger models (70B+) benefit as much from latent anchoring, or if their internal reasoning capabilities render the anchor redundant. Preliminary results suggest that even capable models benefit from the entropy contraction effect, particularly for rigid formatting tasks.

\section{Related Work}
\textbf{Chain-of-Thought (CoT):} CoT relies on the model generating its own reasoning steps in natural language. TSCE subsumes this by treating the ``thought'' as a latent variable (the anchor) that can be optimized directly, rather than relying on the model's default propensity to reason.

\textbf{RLHF / PPO:} Standard RLHF fine-tunes the entire model, which is expensive and can degrade general capabilities (alignment tax). TSCE freezes the base model and only learns a tiny external policy, preserving the base model's knowledge while steering its behavior.

\textbf{Soft Prompting:} Continuous soft prompts (e.g., Prefix-Tuning) require access to the model's embedding layer and cannot be used with black-box APIs. TSCE's discrete anchors are plain text tokens, compatible with any LLM endpoint.

\section{Conclusion}
Two-Step Contextual Enrichment (TSCE) provides a robust framework for controlling Large Language Models through learned, discrete latent anchors. By combining a diffusion-style anchor policy with a rigorous RL loop and intrinsic reasoning rewards (the Thought Panel), TSCE achieves significant performance gains on complex tasks. It offers a scalable, model-agnostic path toward reliable AI systems that can reason and adhere to strict constraints without the need for expensive fine-tuning.

\end{document}
